%% mathstc-binomiale-batons — loi binomiale B(20 ; 0,1) : diagramme en bâtons de P(X = k).
%% La valeur la plus probable se lit sur le bâton le plus haut ; l'espérance E(X) = np = 2
%% est le « centre de gravité » de la distribution. La loi est dissymétrique (queue à
%% droite) : elle ne serait symétrique que pour p = 0,5.
%% Échelles : 1 rang -> 0,7 cm ; 1 unité de probabilité -> 15 cm.
\documentclass[tikz,border=4pt]{standalone}
\usepackage{liaisons}
\begin{document}
\begin{tikzpicture}[x=0.7cm,y=15cm,
  axe/.style={-{Stealth[length=5pt]}, line width=0.6pt},
  baton/.style={solideA, line width=3pt},
  grille/.style={line width=0.3pt, black!15},
  grad/.style={font=\scriptsize, inner sep=2.5pt},
  note/.style={font=\scriptsize, text=black!60, inner sep=1.5pt}]

%% ================= quadrillage horizontal =================================
\foreach \p in {0.05,0.10,0.15,0.20,0.25} { \draw[grille] (-0.5,\p) -- (8.7,\p); }

%% ================= axes et graduations ====================================
\draw[axe] (-0.6,0) -- (9.2,0) node[anchor=north east, inner sep=4pt, font=\small] {$k$};
\draw[axe] (-0.5,-0.012) -- (-0.5,0.315)
     node[anchor=south west, inner sep=2pt, font=\small] {$P(X=k)$};
\foreach \k in {0,1,...,8} {
  \draw[line width=0.5pt] (\k,0) -- (\k,-0.006);
  \node[grad, anchor=north] at (\k,-0.006) {\k}; }
\foreach \p/\lab in {0.05/{0{,}05}, 0.10/{0{,}10}, 0.15/{0{,}15}, 0.20/{0{,}20}, 0.25/{0{,}25}} {
  \draw[line width=0.5pt] (-0.5,\p) -- (-0.62,\p);
  \node[grad, anchor=east] at (-0.62,\p) {$\lab$}; }

%% ================= les neuf bâtons ========================================
\foreach \k/\p in {0/0.1216, 1/0.2702, 2/0.2852, 3/0.1901, 4/0.0898,
                   5/0.0319, 6/0.0089, 7/0.0020, 8/0.0004} {
  \draw[baton] (\k,0) -- (\k,\p);
  \fill[solideA] (\k,\p) circle (1.8pt); }

%% ================= quatre valeurs seulement ===============================
\foreach \k/\lab/\p in {0/{0{,}122}/0.1216, 1/{0{,}270}/0.2702,
                        2/{0{,}285}/0.2852, 3/{0{,}190}/0.1901} {
  \node[grad, anchor=south] at (\k,\p) {$\lab$}; }

%% ================= espérance ==============================================
\draw[-{Stealth[length=5pt]}, solideE, line width=1.1pt] (2,-0.030) -- (2,-0.008);
\node[grad, text=solideE, anchor=north] at (2,-0.032) {$E(X)=np=2$};

%% ================= loi et dissymétrie =====================================
\node[draw=black!45, line width=0.5pt, rounded corners=2pt, fill=white, inner sep=3.5pt,
      font=\small, anchor=north east] at (8.6,0.295) {$\mathcal{B}(20\,;0{,}1)$};
\node[note, anchor=north east] at (8.7,-0.062) {la loi n'est symétrique que si $p=0{,}5$};
\end{tikzpicture}
\end{document}
